\section{Ethics And Scope}

\eha{}-Uncued is a diagnostic benchmark for controlled polluted evidence environments. It is not a safety certification, a deployment approval test, or a general leaderboard. High performance on this pilot would not establish that an agent is reliable in real legal, medical, scientific, or financial workflows.

The benchmark uses synthetic entities and synthetic evidence packets. This avoids collecting sensitive real-user data and makes gold provenance explicit. The tradeoff is ecological validity: synthetic documents cannot cover the full distribution of open-web evidence, institutional records, or adversarial manipulation.

The intended use is research analysis of failure modes: belief formation, evidence-role hygiene, shortcut sensitivity, and verification-action expression. The intended non-use is automated certification of deployed systems without additional task coverage, independent audits, domain-specific validation, and human oversight.
